% ============================================================
\chapter{Statistical Mechanics: Multiplicity and Entropy}
% ============================================================

% ============================================================
\section{Microstates, Macrostates, and Multiplicity}
% ============================================================

Irreversible processes are not logically forbidden by the laws of mechanics --- they are simply overwhelmingly probable. To understand why, we must count the number of ways a system can be arranged.

Imagine flipping three distinguishable coins. There are $2^3 = 8$ possible outcomes, and these naturally organise into two levels of description.

\begin{definition}
A \textbf{microstate} is a complete specification of the system --- for example, exactly which coins show heads and which show tails. A \textbf{macrostate} is a coarser description, specifying only a macroscopic quantity such as the total number of heads. The \textbf{multiplicity} $\Omega$ of a macrostate is the number of microstates consistent with it.
\end{definition}

For $N$ distinguishable objects of which $n$ are in a designated state, the multiplicity is the binomial coefficient
\[
    \Omega(N, n) = \binom{N}{n} = \frac{N!}{n!\,(N-n)!}.
\]
Systems evolve toward macrostates of higher multiplicity because there are simply more ways to be in them. The macrostate with one head out of three coins ($\Omega = 3$) is three times more likely than the macrostate with zero heads ($\Omega = 1$), and for large $N$ the disparity becomes astronomically pronounced.

\begin{definition}[Fundamental Assumption of Statistical Mechanics]
In an isolated system in thermal equilibrium, every accessible microstate is equally probable.
\end{definition}

This assumption implies that any microscopic process $X \to Y$ and its time-reverse $Y \to X$ occur at the same rate. It does \emph{not} mean all macrostates are equally likely: a macrostate with more microstates is simply more probable. For the assumption to be self-consistent, the system must have some mechanism --- however weak --- to explore all accessible microstates over time.

\subsection{Two-State Paramagnet}

A paramagnet is a material whose constituent particles behave like tiny magnetic compass needles that tend to align parallel to an applied field, but only while the field is present (unlike a ferromagnet, which maintains alignment spontaneously). Each particle carries a magnetic \textbf{dipole} moment that may belong to a whole group of atoms, a single atom, or a single electron.

In the idealised \textbf{two-state paramagnet}, each dipole is constrained to point either parallel ($\uparrow$) or antiparallel ($\downarrow$) to an applied field. With $N$ total dipoles and $N_\uparrow$ pointing up, the number of down-spins is $N_\downarrow = N - N_\uparrow$, and the multiplicity of a configuration with $N_\uparrow$ up-spins is
\[
    \Omega(N, N_\uparrow) = \binom{N}{N_\uparrow} = \frac{N!}{N_\uparrow!\, N_\downarrow!}.
\]
Parallel dipoles carry less (more negative) energy than antiparallel ones, because flipping a dipole against an external field requires work. Specifying $N_\uparrow$ is therefore sufficient to determine both the energy and the multiplicity of the macrostate.

% ============================================================
\section{The Einstein Model of Solids}
% ============================================================

Consider a solid modelled as a collection of quantum harmonic oscillators whose energies come in discrete units of size $\varepsilon = hf$, where $h = 6.63\times10^{-34}$~J$\cdot$s is Planck's constant and $f$ is the natural oscillation frequency. Classically the potential energy is the continuous function $\tfrac{1}{2}k_s x^2$, but quantum mechanics forces the allowed energies to be integer multiples of $hf$ above the zero-point energy $E_0 = \tfrac{1}{2}hf$.

\begin{center}
[DIAGRAM: Several identical harmonic-oscillator potential wells drawn side by side, with horizontal lines marking equally-spaced energy levels $E_0 = \tfrac{1}{2}hf$, $\tfrac{3}{2}hf$, $\tfrac{5}{2}hf$, \ldots\ inside each well, labeled ``Oscillators''.]
\end{center}

In a three-dimensional solid, each atom sits at a lattice site and vibrates independently along three orthogonal directions, giving $N_{\text{atoms}}$ atoms a total of $N = 3N_{\text{atoms}}$ independent oscillators. The central combinatorial question is: in how many ways can $q$ indistinguishable energy units be distributed among $N$ oscillators?

This is a ``stars-and-bars'' problem. We want to distribute $q$ identical objects (stars) into $N$ bins separated by $N-1$ dividers (bars). The answer involves choosing the positions of the $N-1$ bars among $q + N - 1$ total slots.

\begin{formula}[Multiplicity of an Einstein Solid]
The number of microstates consistent with having exactly $q$ energy units shared among $N$ oscillators is
\[
    \Omega(N,q) = \binom{q+N-1}{q} = \frac{(q+N-1)!}{q!\,(N-1)!}.
\]
\end{formula}

As a quick check: the arrangement $*\,|\,{*}{*}\,|\,|\,{*}{*}$ uses 3 bars and describes 4 oscillators carrying 5 energy units, consistent with $\Omega(4,5) = \binom{8}{5} = 56$.

\subsection{Interacting Einstein Solids and Equilibrium}

Consider two Einstein solids, A and B, placed in thermal contact. They are \emph{weakly coupled}: energy can flow between them, but the rate of inter-solid exchange is much slower than intra-solid redistribution. Consequently, on short timescales the individual energies $U_A$ and $U_B$ are essentially fixed, while on long timescales only the total $U_{\text{tot}} = U_A + U_B$ is conserved. The macrostate of the combined system is specified by the pair $(q_A, q_B)$ with $q_A + q_B = q$ constant.

\begin{center}
[DIAGRAM: Two rectangular boxes labeled ``Solid A, $N_A$, $q_A$'' and ``Solid B, $N_B$, $q_B$'' connected by a wavy arrow labeled ``Energy'', enclosed in a dashed rectangle labeled ``Isolated system''.]
\end{center}

Each solid has its own multiplicity:
\[
    \Omega_A = \binom{q_A + N_A - 1}{q_A}, \qquad
    \Omega_B = \binom{q_B + N_B - 1}{q_B}.
\]
Because the two solids are statistically independent, the total multiplicity factorises:
\[
    \Omega_{\text{tot}} = \Omega_A \cdot \Omega_B.
\]
The distribution of $\Omega_{\text{tot}}$ over macrostates peaks sharply near $q_A = q_B = q/2$ for equal solids. If the system starts in an extreme macrostate with nearly all energy in one solid, it is overwhelmingly likely to evolve toward the peak. This is the statistical content of the second law: energy (heat) flows spontaneously from the hotter solid to the cooler one, because the macrostates near equilibrium carry the overwhelming majority of the total multiplicity.

% ============================================================
\section{Large Systems and Stirling's Approximation}
% ============================================================

As $N$ and $q$ grow, the multiplicity distribution $\Omega_{\text{tot}}(q_A)$ becomes dramatically sharper relative to the total range of $q_A$. Two mathematical observations underlie everything that follows. First, when $\lambda \gg a$, we have $\lambda + a \approx \lambda$ (e.g., $10^{23} + 23 \approx 10^{23}$). Second, the Taylor expansion $\ln(1+x) \approx x$ for $|x| \ll 1$. Together with a large-$N$ approximation for factorials, these yield all the tools we need.

\begin{formula}[Stirling's Approximation]
For $N \gg 1$,
\[
    N! \approx N^N e^{-N} \sqrt{2\pi N},
    \qquad \text{equivalently,} \qquad
    \ln(N!) \approx N\ln N - N.
\]
The logarithmic form follows from approximating the sum $\ln N! = \sum_{k=1}^N \ln k$ by the integral $\int_1^N \ln x\,dx = [x\ln x - x]_1^N \approx N\ln N - N$. The subleading term $\tfrac{1}{2}\ln(2\pi N)$ arises from a saddle-point correction and is significant when one needs the prefactor.
\end{formula}

\subsection{High-Temperature Limit ($q \gg N$)}

When many more energy units are present than oscillators, each oscillator typically carries many quanta. In this regime we may approximate $(q + N - 1)! \approx (q+N)!$ and apply Stirling's approximation to $\ln\Omega$. Since $q \gg N$, we expand $\ln(q + N) = \ln q + \ln(1 + N/q) \approx \ln q + N/q$ and find, after collecting terms:
\begin{align*}
    \ln\Omega
    &= (q+N)\ln(q+N) - q\ln q - N\ln N + \cdots \\
    &\approx N\ln\!\left(\frac{q}{N}\right) + N.
\end{align*}

\begin{formula}[Einstein Solid: High-Temperature Limit]
For $q \gg N$,
\[
    \Omega(N,q) \approx \left(\frac{eq}{N}\right)^N.
\]
\end{formula}

\subsection{Low-Temperature Limit ($q \ll N$)}

When few energy units are present, most oscillators sit in their ground state. Expanding $\ln(N + q) \approx \ln N + q/N$ and repeating the same steps yields
\[
    \Omega(N,q) \approx \left(\frac{eN}{q}\right)^q \qquad (q \ll N).
\]
In both limits, and for the general case in between, the key feature is that $\Omega$ depends on $N$ and $q$ through the ratio $q/N$ --- the number of energy units per oscillator, which is directly proportional to the temperature.

\subsection{Sharpness of the Multiplicity Peak}

Consider two large, interacting Einstein solids each with $N$ oscillators sharing $q$ total energy units in the high-temperature limit $q \gg N$. Using the high-temperature formula for each solid:
\[
    \Omega_{\text{tot}} = \left(\frac{e}{N}\right)^{2N}(q_A q_B)^N.
\]
This is maximised when $q_A q_B$ is maximised subject to $q_A + q_B = q$, which occurs at $q_A = q_B = q/2$:
\[
    \Omega_{\max} = \left(\frac{eq}{2N}\right)^{2N}.
\]
To quantify how sharply $\Omega_{\text{tot}}$ peaks, write $q_A = q/2 + x$ and $q_B = q/2 - x$ for small displacement $x$, so $q_A q_B = (q/2)^2 - x^2$. Using $\ln(1 - \epsilon) \approx -\epsilon$ for small $\epsilon = (2x/q)^2$:

\begin{formula}[Gaussian Form of the Multiplicity Peak]
Near the maximum,
\[
    \Omega_{\text{tot}}(x) = \Omega_{\max}\, \exp\!\left[-N\!\left(\frac{2x}{q}\right)^2\right].
\]
The distribution is a Gaussian in the displacement $x = q_A - q/2$ from the equilibrium energy partition.
\end{formula}

The $1/e$ half-width of the Gaussian follows from setting the exponent equal to $-1$:
\[
    N\!\left(\frac{2x}{q}\right)^2 = 1
    \;\Longrightarrow\; \Delta x = \frac{q}{\sqrt{N}}.
\]
The relative width of the peak compared to the full range $q$ is therefore $\Delta x / q = 1/\sqrt{N}$. For a macroscopic solid $N \sim 10^{23}$, this ratio is $\sim 10^{-11.5}$ --- the peak is vanishingly narrow. In the \textbf{thermodynamic limit} ($N \to \infty$), random fluctuations away from the most probable macrostate are so improbable that they are unmeasurable in practice. The second law of thermodynamics is not a fundamental constraint on the possible microstates; it is a statement about probability in the thermodynamic limit.

% ============================================================
\section{Density of States and Semi-Classical Phase Space}
% ============================================================

The multiplicity of a discrete system like the Einstein solid is a straightforward combinatorial count. For a continuous system like an ideal gas, we need a more careful procedure, since classical phase space is a continuum. The semi-classical approach bridges quantum mechanics and classical mechanics.

\subsection{State Counting and One Quantum State Per $h^3$}

For a single particle confined to a cubic box of side $L$, solving the Schr\"{o}dinger equation gives energy eigenvalues
\[
    E(n_1, n_2, n_3) = \frac{\pi^2\hbar^2}{2mL^2}(n_1^2 + n_2^2 + n_3^2),
\]
where $n_1, n_2, n_3 = 1, 2, 3, \ldots$ are positive integers. The cumulative number of states $N(E)$ up to energy $E$ counts the positive-integer lattice points inside a sphere of radius $C = \sqrt{2mL^2 E / (\pi^2\hbar^2)}$. Since the quantum level spacing ($\sim \pi^2\hbar^2/(2mL^2) \approx k_B \times 10^{-13}$~K for helium in a millimeter box) is many orders of magnitude smaller than $k_B T$ at any accessible temperature, we may replace the lattice sum by an integral over the positive octant of this sphere:
\[
    N(E) = \frac{1}{8} \cdot \frac{4\pi}{3} C^3 = \frac{1}{6\pi^2}\left(\frac{2m}{\hbar^2}\right)^{3/2} V E^{3/2}.
\]
where $V = L^3$. The density of states follows by differentiation:
\begin{formula}
\[
    \rho(E) \equiv \dv{N}{E} = \frac{1}{4\pi^2}\left(\frac{2m}{\hbar^2}\right)^{3/2} V\sqrt{E}.
\]
\end{formula}
The density of states grows as $\sqrt{E}$ for a particle in three dimensions.

\subsection{Semi-Classical Phase Space and the $h^3$ Cell}

It is often more convenient to count states using the classical energy $E(\vb{x}, \vb{p})$ rather than quantum numbers. For a single particle, the phase space is six-dimensional: $\{x, y, z, p_x, p_y, p_z\}$. The classical integral over phase space is infinite (a continuum), but quantum mechanics provides a natural coarse-graining: each quantum state occupies a minimum volume in phase space, which we can determine by matching the classical and quantum counts.

The classical cumulative count up to energy $E$ is
\[
    N(E) = \frac{1}{\Delta}\int d^3x\int d^3p\;\Theta\!\bigl(E - E(\vb{x}, \vb{p})\bigr),
\]
where $\Delta$ is the phase-space volume per quantum state. For a particle in a box with $E = p^2/(2m)$, the spatial integral gives $V$ and the momentum integral gives the volume of the sphere $p^2 \leq 2mE$:
\[
    N(E) = \frac{V}{\Delta}\cdot\frac{4\pi}{3}(2mE)^{3/2}.
\]
Comparing to the quantum lattice result requires:
\begin{formula}[Phase Space Volume Per Quantum State]
\[
    \Delta = (2\pi\hbar)^3 = h^3.
\]
\end{formula}
This follows directly from the Heisenberg uncertainty principle: $\Delta x\,\Delta p_x \sim \hbar$ per spatial direction, so a minimum-uncertainty cell in phase space has volume $(2\pi\hbar)^3$ in three dimensions.

\subsection{General $N$-Particle System}

For an $N$-particle system with classical Hamiltonian $H(\vb{x}_1, \ldots, \vb{x}_N, \vb{p}_1, \ldots, \vb{p}_N)$, each quantum state occupies a volume $h^{3N}$ in $6N$-dimensional phase space. The total number of states with energy at most $U$ is:
\begin{formula}
\[
    N(U) = \frac{1}{(2\pi\hbar)^{3N}}\int d^{3N}x\int d^{3N}p\;\Theta\!\bigl(U - H(\vb{x}, \vb{p})\bigr).
\]
\end{formula}
For identical particles, particle permutations produce physically indistinguishable configurations, so we must divide by $N!$ to avoid over-counting.

% ============================================================
\section{The Multiplicity of an Ideal Gas: Sackur-Tetrode}
% ============================================================

We now apply the semi-classical framework to compute the multiplicity and entropy of an ideal gas of $N$ identical monatomic particles in a box of volume $V$ with total energy $U$.

\subsection{Volume of the $3N$-Dimensional Momentum Sphere}

For a non-interacting gas, the Hamiltonian is $H = \sum_{j=1}^N \vb{p}_j^2/(2m)$, and the spatial integral over $d^{3N}x$ gives $V^N$. The momentum constraint $\sum_j \vb{p}_j^2 \leq 2mU$ restricts the $3N$-dimensional momentum vector to a ball of radius $\sqrt{2mU}$. The volume of the unit sphere in $d = 3N$ dimensions is
\[
    V_{3N} = \frac{\pi^{3N/2}}{\Gamma\!\left(\frac{3N}{2} + 1\right)},
\]
where the Gamma function satisfies $\Gamma(n) = (n-1)!$ for positive integers and $\Gamma(n+1) = n\,\Gamma(n)$ in general. Putting it together and dividing by $h^{3N}$ (state counting) and $N!$ (indistinguishability):
\[
    \Omega(U, V, N) \approx \frac{1}{N!}\cdot\frac{V^N}{h^{3N}}\cdot\frac{\pi^{3N/2}}{\left(\tfrac{3N}{2}\right)!}\,(2mU)^{3N/2},
\]
where for large $3N$ we use $\Gamma\!\left(\tfrac{3N}{2}+1\right) \approx \left(\tfrac{3N}{2}\right)!$.

\subsection{Derivation of the Sackur-Tetrode Equation}

To find $S = k\ln\Omega$, we apply Stirling's approximation to both $\ln N!$ and $\ln\!\left(\tfrac{3N}{2}\right)!$. Writing $\ln N! \approx N\ln N - N$ and $\ln\!\left(\tfrac{3N}{2}\right)! \approx \tfrac{3N}{2}\ln\tfrac{3N}{2} - \tfrac{3N}{2}$ and collecting all terms:
\begin{align*}
    \ln\Omega
    &= N\!\left[-\ln N + 1 + \ln V - 3\ln h + \frac{3}{2}\ln\pi + \frac{3}{2}\ln(2mU)
       - \frac{3}{2}\ln\frac{3N}{2} + \frac{3}{2}\right] \\
    &= N\!\left[\ln\frac{V}{N} + \frac{3}{2}\ln\!\left(\frac{4\pi mU}{3Nh^2}\right) + \frac{5}{2}\right].
\end{align*}

\begin{formula}[Sackur-Tetrode Equation]
\[
    \boxed{S(N,V,U) = Nk\!\left[\ln\!\left(\frac{V}{N}\left(\frac{4\pi m U}{3Nh^2}\right)^{3/2}\right) + \frac{5}{2}\right].}
\]
\end{formula}

Writing $c_0 = \tfrac{3}{2}\ln\!\left(\tfrac{4\pi m}{3h^2}\right) + \tfrac{5}{2}$ as a species-dependent constant, the entropy takes the compact form
\[
    S = Nk\!\left[\ln\frac{V}{N} + \frac{3}{2}\ln\frac{U}{N} + c_0\right].
\]
The entropy increases with $N$, $V$, and $U$: more particles, more volume, or more energy each mean more accessible microstates.

\subsection{The Gibbs Paradox and the $1/N!$ Factor}

The $1/N!$ factor in the multiplicity is essential for the entropy to be \emph{extensive} (scaling linearly with system size). To see what goes wrong without it, consider doubling a gas from $(U, V, N)$ to $(2U, 2V, 2N)$. Without the $1/N!$ factor the entropy would scale as
\[
    S(2U, 2V, 2N) = 2S(U, V, N) + 2kN\ln 2 \neq 2S(U,V,N).
\]
The spurious extra term $2kN\ln 2$ is the \textbf{Gibbs paradox}: an apparent entropy increase even when two portions of the \emph{same} gas are combined. This would imply that simply relabelling a partition position could create entropy, which is physically absurd.

The resolution is that identical particles are genuinely indistinguishable. Swapping two helium atoms produces no new microstate, and the $1/N!$ corrects for this over-counting. With the correct factor, one can verify directly that the Sackur-Tetrode entropy is exactly extensive: substituting $(2U, 2V, 2N)$ gives $2S(U, V, N)$ with no remainder.

% ============================================================
\section{Entropy: $S = k\ln\Omega$}
% ============================================================

The pattern across every system we have studied is that any large system will almost certainly be found in or very near the macrostate of greatest multiplicity. Fluctuations away from this peak are not forbidden, only exponentially suppressed. This observation is the statistical underpinning of the second law.

\begin{definition}[Boltzmann's Entropy]
The \textbf{entropy} of a macrostate with multiplicity $\Omega$ is
\[
    S \equiv k\ln\Omega,
\]
with $k = 1.38\times10^{-23}$~J/K Boltzmann's constant. Since $\Omega \geq 1$, entropy is non-negative. Since $\Omega_{\text{tot}} = \Omega_A\Omega_B$ for independent subsystems, entropy is \textbf{additive}: $S_{\text{tot}} = S_A + S_B$.
\end{definition}

\begin{theorem}[Second Law of Thermodynamics]
The total entropy of an isolated system tends to increase (or remain constant) over time. Equivalently, the multiplicity tends to increase, because a randomly chosen microstate is overwhelmingly likely to belong to a high-multiplicity macrostate.
\end{theorem}

The connection to the thermodynamic (Clausius) entropy $dS = \dbar Q/T$ will be established in Section~\ref{sec:temperature-from-entropy}, where the statistical definition is shown to reproduce the correct ideal gas temperature and pressure, making the identification exact.

\subsection{Entropy of Mixing}

When two \emph{different} ideal gases, initially separated, are allowed to mix, each gas expands into the combined volume and the total entropy increases. If gas $A$ contains $(1-x)N$ molecules and gas $B$ contains $xN$ molecules, each initially at the same temperature and pressure in their respective partial volumes $(1-x)V$ and $xV$, then after mixing each species expands to occupy the full volume $V$:
\[
    \Delta S_A = (1-x)Nk\ln\frac{1}{1-x}, \qquad \Delta S_B = xNk\ln\frac{1}{x}.
\]

\begin{formula}[Entropy of Mixing]
\[
    \Delta S_{\text{mix}} = -Nk\bigl[(1-x)\ln(1-x) + x\ln x\bigr].
\]
\end{formula}

This result is non-negative for all $0 \leq x \leq 1$ (since both logarithm terms are non-positive), and vanishes only for pure species ($x = 0$ or $x = 1$). The combinatorial derivation is equally illuminating: we can derive this directly by noting that $\Delta S_{\text{mix}} = k\ln\binom{N}{N_A}$ counts the number of distinguishable ways to interleave $(1-x)N$ particles of species $A$ with $xN$ particles of species $B$, and applying Stirling's approximation to the binomial coefficient.

\begin{center}
[DIAGRAM: Two separate boxes (gas $A$ left, gas $B$ right, both at same $T$ and $P$) with an arrow to a single combined box containing both gases mixed.]
\end{center}

When the two gases are \emph{identical}, the same calculation gives $\Delta S_{\text{mix}} = 0$: combining two equal portions of the same gas at the same temperature and pressure creates no new entropy. This is guaranteed by the $1/N!$ factor in the Sackur-Tetrode multiplicity, which correctly encodes the indistinguishability of identical particles. Swapping two identical molecules produces no new microstate, so the system is already at its maximum entropy and no further equilibration is required.

\subsection{Free Expansion}

\begin{center}
[DIAGRAM: Gas confined to the left half of an insulated box by a removable partition; arrow; gas fills the entire box after partition removal. The walls are insulated ($Q = 0$) and the right side is vacuum ($W = 0$).]
\end{center}

In a free expansion, a partition is removed and the gas expands into vacuum. No work is done ($W = 0$) and no heat is transferred ($Q = 0$), so the energy $U$ is unchanged. Yet the volume doubles, and the Sackur-Tetrode equation gives:

\begin{formula}
\[
    \Delta S = Nk\ln\frac{V_f}{V_i} > 0 \quad \text{(irreversible).}
\]
\end{formula}

This entropy increase is genuine: the process is irreversible, and entropy has been created. Unlike a quasistatic expansion against a piston --- which is reversible and for which $\Delta S = \dbar Q/T$ applies --- in free expansion there is no external reservoir to absorb entropy, and the process cannot be run in reverse without violating the second law.

% ============================================================
\section{Temperature and Pressure from Entropy}
\label{sec:temperature-from-entropy}
% ============================================================

We have defined entropy $S = k\ln\Omega$ and seen that an isolated system tends toward its highest-entropy macrostate. We now derive how entropy relates to the familiar thermodynamic variables temperature and pressure.

\subsection{The Statistical Definition of Temperature}

Two objects are at the same temperature when they are in thermal equilibrium; two objects in thermal equilibrium simultaneously maximise the total entropy. Suppose subsystems $A$ and $B$ are in thermal contact, able to exchange energy $\dbar Q$ but not volume or particles, with $U_A + U_B = U_{\text{tot}}$ fixed. The condition for a maximum of $S_{\text{tot}} = S_A + S_B$ with respect to the energy split is
\[
    \pdv{S_{\text{tot}}}{U_A} = 0 \quad\Longrightarrow\quad \left.\pdv{S_A}{U_A}\right|_{V_A,N_A} = \left.\pdv{S_B}{U_B}\right|_{V_B,N_B},
\]
where we used $dU_A = -dU_B$.

\begin{center}
[DIAGRAM: Graph of $S_A$, $S_B$, and $S_{\text{total}}$ as functions of $q_A$ (energy in solid $A$). $S_A$ increases monotonically, $S_B$ decreases monotonically, $S_{\text{total}}$ has a smooth maximum. The equilibrium point is where the slopes of $S_A$ and $S_B$ are equal in magnitude.]
\end{center}

The intuition is as follows. At a point where $\partial S_A/\partial U_A > \partial S_B/\partial U_B$, transferring a small amount of energy from $B$ to $A$ increases the total entropy: the entropy gained by $A$ exceeds the entropy lost by $B$. Energy therefore flows from $B$ into $A$. A steeper $\partial S/\partial U$ corresponds to a colder object (it gains more entropy per unit energy received), while a shallower slope corresponds to a hotter object. This motivates:

\begin{definition}[Statistical Definition of Temperature]
\[
    \frac{1}{T} \equiv \left(\pdv{S}{U}\right)_{V,N}.
\]
Equivalently, $dS = dU/T$ at constant $N$ and $V$.
\end{definition}

\subsection{Pressure from Entropy}

By an exactly analogous argument, if two systems can exchange volume through a movable piston (at fixed total volume $V_A + V_B$ and fixed energies), the maximum-entropy condition is
\[
    \left.\pdv{S_A}{V_A}\right|_{U_A, N_A} = \left.\pdv{S_B}{V_B}\right|_{U_B, N_B}.
\]
This quantity must equal $P/T$ for consistency with the thermodynamic relation $dU = T\,dS - P\,dV$, giving:
\begin{formula}
\[
    \frac{P}{T} = \left(\pdv{S}{V}\right)_{U,N}.
\]
\end{formula}

\subsection{Verification Against the Ideal Gas}

The Sackur-Tetrode entropy has $U$-dependence $S = \tfrac{3}{2}Nk\ln U + (\text{function of }V, N)$, so
\[
    \frac{1}{T} = \left(\pdv{S}{U}\right)_{V,N} = \frac{3Nk}{2U}
    \qquad\Longrightarrow\qquad \boxed{U = \tfrac{3}{2}NkT.}
\]
This recovers the ideal gas internal energy law purely from counting microstates. The $V$-dependence of the Sackur-Tetrode entropy is $S = Nk\ln V + \cdots$, so
\[
    \frac{P}{T} = \left(\pdv{S}{V}\right)_{U,N} = \frac{Nk}{V}
    \qquad\Longrightarrow\qquad \boxed{PV = NkT.}
\]
The ideal gas law emerges as a consequence of counting microstates in phase space, with no additional assumptions. The definition $1/T = \partial S/\partial U$ is not circular: it is defined purely in terms of the entropy function $S(U,V,N)$ and reproduces the correct thermodynamics without free parameters.

\subsection{Connecting to Clausius: $dS = \dbar Q/T$}

When a small amount of heat $\dbar Q$ is added to a system at constant volume (so no work is done), the first law gives $dU = \dbar Q$, and the temperature definition gives $dS = dU/T = \dbar Q/T$. This is exactly the Clausius entropy. For a reversible process one can integrate:
\begin{formula}
\[
    \Delta S = \int_{T_i}^{T_f}\frac{C_V}{T}\,dT,
\]
\end{formula}
which allows entropy to be measured calorimetrically from $C_V(T)$ alone, without any microscopic model.

\subsection{The Zeroth Law as a Corollary}

The statistical definition immediately implies the Zeroth Law. If $A$ is in thermal equilibrium with $C$, both have the same $\partial S/\partial U = 1/T$. If $B$ is also in equilibrium with $C$, then $B$ shares the same $\partial S/\partial U$ as $C$, and hence the same as $A$. Therefore $A$ and $B$ are in mutual equilibrium.

\begin{theorem}[Zeroth Law of Thermodynamics]
If system $A$ is in thermal equilibrium with system $C$, and system $B$ is in thermal equilibrium with system $C$, then $A$ is in thermal equilibrium with $B$. This follows directly from the entropy-maximisation definition of temperature: $(\partial S_A/\partial U_A)_{V,N} = (\partial S_C/\partial U_C)_{V,N}$ and $(\partial S_B/\partial U_B)_{V,N} = (\partial S_C/\partial U_C)_{V,N}$ together imply $(\partial S_A/\partial U_A)_{V,N} = (\partial S_B/\partial U_B)_{V,N}$, so temperature is a transitive, equivalence-class-defining property of matter.
\end{theorem}

The Third Law also emerges naturally: at $T \to 0$, the system settles into its unique quantum ground state, so $\Omega \to 1$ and $S \to k\ln 1 = 0$.

% ============================================================
\section{Paramagnetism}
% ============================================================

The two-state paramagnet is one of the simplest systems for which we can carry the microcanonical programme to completion, computing $U(T)$, $M(T)$, and $S(T)$ analytically from $S = k\ln\Omega$.

\subsection{Setup and Multiplicity}

We study a system of $N$ non-interacting spin-$\tfrac{1}{2}$ dipoles immersed in a constant external magnetic field $B$ in the $+\hat{z}$ direction. Each dipole can point either ``up'' ($+\hat{z}$, aligned with $B$, lower energy) or ``down'' ($-\hat{z}$, anti-aligned, higher energy), with energies
\[
    \varepsilon_{\uparrow} = -\mu B \quad (\text{aligned, lower energy}), \qquad
    \varepsilon_{\downarrow} = +\mu B \quad (\text{anti-aligned, higher energy}),
\]
where $\mu > 0$ is the magnitude of the dipole's magnetic moment. With $N_\uparrow$ up-spins and $N_\downarrow = N - N_\uparrow$ down-spins, the total energy is
\[
    U = N_\uparrow(-\mu B) + N_\downarrow(\mu B) = \mu B(N - 2N_\uparrow).
\]

\begin{definition}
The \textbf{magnetization} $M$ is the total magnetic moment of the system:
\[
    M \equiv \mu(N_\uparrow - N_\downarrow) = -\frac{U}{B}.
\]
\end{definition}

The multiplicity depends only on $N_\uparrow$ (equivalently on the energy or magnetization):
\begin{formula}
\[
    \Omega(N_\uparrow) = \binom{N}{N_\uparrow} = \frac{N!}{N_\uparrow!\,N_\downarrow!}.
\]
\end{formula}

\subsection{Temperature, Energy, and Magnetization}

To find $U(T)$, we differentiate the entropy $S = k\ln\Omega$ with respect to $U$. Using Stirling's approximation,
\[
    S \approx k\!\left[N\ln N - N_\uparrow\ln N_\uparrow - N_\downarrow\ln N_\downarrow\right],
\]
where $N_\uparrow = (N - U/\mu B)/2$ and $N_\downarrow = (N + U/\mu B)/2$. Differentiating:
\[
    \frac{1}{T} = \pdv{S}{U} = \frac{k}{2\mu B}\ln\frac{N_\downarrow}{N_\uparrow}.
\]
Solving for the ratio $N_\uparrow / N_\downarrow = e^{2\mu B/kT}$ and using $N_\uparrow + N_\downarrow = N$, we find $N_\uparrow - N_\downarrow = N\tanh(\mu B/kT)$. Therefore:

\begin{formula}
\[
    U(T) = -N\mu B\tanh\!\left(\frac{\mu B}{kT}\right), \qquad
    M(T) = N\mu\tanh\!\left(\frac{\mu B}{kT}\right).
\]
\end{formula}

\subsection{Limiting Behaviour and the Saturation of Magnetization}

At low temperature ($kT \ll \mu B$), the $\tanh$ saturates to unity: $U \to -N\mu B$ (all spins aligned in the ground state) and $M \to N\mu$ (full saturation magnetization). At high temperature ($kT \gg \mu B$), $\tanh(\mu B/kT) \approx \mu B/kT$, giving
\[
    M \approx \frac{N\mu^2 B}{kT} \quad (kT \gg \mu B).
\]
This is the \textbf{Curie law}: magnetization is proportional to $B/T$ and falls off as temperature increases, because thermal fluctuations increasingly disorder the spins.

\begin{center}
[DIAGRAM: Plot of $M(T)/(N\mu)$ vs.\ $kT/(\mu B)$. The curve starts at 1 for $T = 0$ (saturation), decreases monotonically, and approaches zero as $T \to \infty$ following the Curie law $M \propto 1/T$ (shown as a dashed asymptote).]
\end{center}

\subsection{Entropy of the Paramagnet}

Substituting $N_\uparrow/N = [1 + \tanh(\mu B/kT)]/2$ back into $S = k\ln\Omega$:
\[
    \frac{S}{Nk} = -\frac{1+t}{2}\ln\frac{1+t}{2} - \frac{1-t}{2}\ln\frac{1-t}{2}, \quad t \equiv \tanh\!\left(\frac{\mu B}{kT}\right).
\]
At $T = 0$ all spins are in the ground state ($N_\uparrow = N$, $\Omega = 1$), so $S = 0$ --- consistent with the Third Law. At $T \to \infty$ the two spin states become equally probable ($N_\uparrow = N_\downarrow = N/2$, $\Omega = \binom{N}{N/2} \approx 2^N$), and $S \to Nk\ln 2$. The entropy rises monotonically from zero to this maximum as thermal fluctuations progressively disorder the spin configuration, reflecting the growth of accessible microstates with increasing temperature.
